\chapter{Estado del Proyecto y Roadmap}
\label{ch:roadmap}
% ==============================================================

\section{Estado actual (agosto 2026)}

El proyecto se planteó originalmente con una duración de 6 meses (enero--junio 2026,
capítulo~\ref{ch:overview}); en la práctica el desarrollo se extendió hasta agosto de 2026 para
completar Auto Mode, el instalador de escritorio y una ronda de pulido orientada a producción
que no estaba en el alcance original.

\begin{table}[H]
\centering
\caption{Estado de hitos del proyecto}
\label{tab:roadmap}
\begin{tabular}{C{1cm} L{5cm} C{1.5cm} L{5cm}}
\toprule
\textbf{Mes} & \textbf{Objetivo} & \textbf{Estado} & \textbf{Nota} \\
\midrule
1 & Preparación datos + diseño interfaces & \textcolor{OliveGreen}{Hecho} & Completado \\
2 & Calibración geométrica                & \textcolor{OliveGreen}{Hecho} & 6 perfiles generados \\
3 & Segmentación (prototipo)              & \textcolor{OliveGreen}{Hecho} & SAM funcional \\
4 & Extracción línea + georreferenciación & \textcolor{OliveGreen}{Hecho} & GeoJSON en EPSG:25830 \\
5 & Validación y robustez                 & \textcolor{OliveGreen}{Hecho} & Las 6 cámaras calibradas y validadas \\
6 & Integración y entrega                 & \textcolor{OliveGreen}{Hecho} & Instalador de escritorio generado y probado \\
\bottomrule
\end{tabular}
\end{table}

\section{Completado en la iteración de junio 2026}

\begin{itemize}[leftmargin=1.5em]
  \item[$\bullet$] Módulo de alineación masiva no-destructiva (\texttt{batch\_alignment.py})
    \begin{itemize}
      \item Router FastAPI \texttt{/api/batch/} con 10 endpoints
      \item Pipeline async SIFT+FLANN+RANSAC sobre lote completo
      \item Mapa de delta compuesto (heatmap HOT + edge-diff cian + métrica Farneback)
      \item Two-phase commit: staging \texttt{/tmp/} → escritura permanente sólo en commit
    \end{itemize}
  \item[$\bullet$] Interfaz de validación por lotes (\texttt{BatchAlignment.vue})
  \item[$\bullet$] Corrección de máscaras SIFT para CAM~1 (problema de gaviotas)
  \item[$\bullet$] Integración en paso 3 del flujo de calibración (\texttt{Calibration.vue})
  \item[$\bullet$] Documentación completa (esta memoria)
\end{itemize}

\section{Fase 3 — Automatización y robustez (completada)}

\begin{enumerate}[leftmargin=1.5em]
  \item \textbf{Procesamiento por lotes}: resuelto por \textbf{Auto Mode}
        (\texttt{backend/auto\_mode.py}, capítulo~\ref{ch:automode}) — descarga, alinea y analiza
        un rango de fechas completo de una cámara de punta a punta, bajo demanda.
  \item \textbf{Filtros de calidad}: \texttt{analyze\_roi()} rechaza automáticamente imágenes con
        confianza de SAM por debajo del umbral configurado (\texttt{conf\_min}), registrando el
        motivo en el log de actividad.
  \item \textbf{Persistencia GeoJSON}: \texttt{\_append\_history()} acumula cada análisis en
        \texttt{coastline\_history\_cam\{id\}.json}, consumido por el gráfico "Tendencia
        Histórica" del Dashboard (\texttt{GET /api/historical-data}).
  \item \textbf{Recalibración de las 6 cámaras}: completada — todas por debajo del objetivo de
        RMSE~$<$~2~m (\texttt{RMSE\_GOOD\_M}, capítulo~\ref{ch:calibration}).
\end{enumerate}

\section{Fase 4 — Entrega e integración (completada)}

\begin{enumerate}[leftmargin=1.5em]
  \item Dashboard de histórico: gráfica de evolución del área seca ya en producción
        (\texttt{Dashboard.vue} + \texttt{GET /api/historical-data}).
  \item Exportación de informes: \texttt{GET /api/report} (CSV/JSON) desde el Dashboard, más un
        informe PDF individual por cámara con el resumen de su calibración
        (\texttt{GET /api/cameras/\{id\}/calibration-report.pdf}, capítulo~\ref{ch:calibration}).
  \item Manual de usuario: \texttt{docs/user\_manual.md}, con la guía paso a paso del asistente
        de calibración de 7 pasos y de Auto Mode.
  \item Instalador de escritorio para Windows: \texttt{installer/cv-lit.iss} +
        \texttt{backend/desktop\_launcher.spec}, compilado y probado
        (\texttt{installer/output/LineaDeCosta-Setup.exe}) — ver capítulo~\ref{ch:setup}.
\end{enumerate}

\section{Ronda de pulido para producción (julio--agosto 2026)}

Tras cerrar las Fases 3 y 4, una auditoría completa del proyecto orientada a llevarlo a
producción añadió, sin estar en el roadmap original:

\begin{itemize}[leftmargin=1.5em]
  \item Endurecimiento de seguridad: \texttt{CORS allow\_credentials=False}, validación de
        \texttt{PUT} en varios endpoints, sanitización de nombres de fichero
        (\texttt{\_safe\_filename()}), límite mínimo de subida (\texttt{MIN\_UPLOAD}), y
        eliminación de credenciales de Obscape hardcodeadas (ahora por variables de entorno,
        \texttt{acces\_api/obscape\_api.py}).
  \item Eliminación de código y scripts no usados (rutas de alineación duplicadas, herramientas
        de calibración de un solo uso, la vista "rectificada UTM" nunca terminada).
  \item Aviso de geometría inestable en el cálculo de homografía (\texttt{geometry\_warning},
        capítulo~\ref{ch:calibration}) y su indicador correspondiente en el stepper del
        asistente.
  \item Suite de tests ampliada: \texttt{tests/test\_auto\_mode.py} y
        \texttt{tests/test\_geometry\_warning.py}, cubriendo Auto Mode y el aviso de geometría
        sin depender de la red real de Obscape.
  \item Esta memoria: reescrita y ampliada de forma sustancial (capítulos~\ref{ch:calibration},
        \ref{ch:batch}, \ref{ch:automode} y auditoría de coherencia del resto de capítulos).
\end{itemize}

% ==============================================================
